\section{Introduction}

Public-sector and enterprise endpoint fleets accumulate disclosed vulnerabilities faster than operations teams can remediate them within constrained maintenance windows. The practical triage question is not \emph{which vulnerabilities are inherently severe} but \emph{which (host, CVE) pairs should receive the next unit of limited remediation capacity} --- a distinction that requires combining exploit-likelihood signals with knowledge of which hosts are most exposed and least well-maintained.

The Exploit Prediction Scoring System (EPSS)~\cite{jacobs2021epss} has emerged as the leading public signal for per-CVE exploit likelihood, outperforming CVSS base scores as a predictor of near-term exploitation in the wild~\cite{jacobs2020improving}. EPSS is incorporated into risk-based vulnerability management (RBVM) tools, CISA vulnerability prioritization guidance~\cite{cisa2021bod2201,cisa2022bod2301}, and remediation policy frameworks. However, EPSS is asset-agnostic by design: it assigns a single probability score to a CVE regardless of whether the vulnerable host is a freshly-telemetered, heavily-patched workstation or an unmanaged endpoint with stale patch data and over-privileged AD group memberships.

Intuitively, a CVE with moderate EPSS on a host that has gone 21 days without telemetry, carries 40\% unpatched CVE debt, and is accessible from domain-privileged accounts represents a materially different risk than the same CVE on a well-maintained host with fresh telemetry and minimal AD exposure. Yet EPSS rankings treat both identically. Host-level \emph{hygiene posture} --- the aggregate state of patch compliance, identity and access management hygiene, and telemetry coverage --- is a natural complement to exploit-likelihood scoring.

Prior work directly relevant to this problem is \textbf{VulnPrio}~\cite{paper1vulnprio} (referred to throughout as Paper~1), which built a reproducible benchmark for (host, CVE) pair prioritization on a synthetic government-shaped fleet (EEHDA) and evaluated 13 strategies combining EPSS, CVSS, KEV membership, asset criticality, and endpoint exposure signals. Paper~1's central null finding was that the proposed multi-factor model was statistically indistinguishable from the EPSS-only baseline: inter-strategy differences fell within seed-to-seed variation, and no CVE-level augmentation strategy materially outperformed EPSS. Paper~1 concluded that this result exposed capacity-driven and base-rate trade-offs and provided a falsifiable benchmark on which calibrated studies could build --- explicitly framing the result as a foundation for future work rather than an endpoint.

Separately, \textbf{HygieneBench}~\cite{paper3hygienebench} (Paper~3) introduced a reproducible benchmark for cyber-hygiene anomaly detection on a synthetic fleet covering Active Directory identity state, endpoint patch posture, vulnerability exposure, and telemetry freshness. HygieneBench's primary finding was that 86.2\% of evaluated (condition, task, method) configurations fail to outperform a rule baseline, but that patch-vulnerability hygiene (Task T5, $+0.210$ average precision over rule for OCSVM) and group-membership drift (Task T2, $+0.185$ AP for temporal z-score) are exceptions where structured ML methods add meaningful signal.

\textbf{The central hypothesis of this paper} is that the signal missing from Paper~1's multi-factor model was not additional CVE-level features but host-level hygiene posture: patch debt, AD privilege exposure, and telemetry freshness. HygieneBench demonstrated that these dimensions carry discriminative structure detectable from the synthetic fleet telemetry. HygienePrio operationalizes this hypothesis by constructing a Hygiene Risk Score (HRS) from HygieneBench-style host telemetry and integrating it into an EPSS-weighted (host, CVE) scorer via an additive combination with an explicit interaction term.

\subsection{Contributions}

This paper makes the following contributions:

\begin{itemize}
  \item \textbf{C1 --- HygienePrio scorer:} A four-component weighted scoring function $S(h, c) = \alpha \times \mathrm{EPSS}(c) + \beta \times \mathrm{HRS}(h) + \gamma \times \mathrm{KEV\_recency}(c) + \delta \times (\mathrm{EPSS}(c) \times \mathrm{HRS}(h))$ that integrates host-level hygiene posture with per-CVE exploit likelihood in a principled, interpretable form.

  \item \textbf{C2 --- Hygiene Risk Score (HRS):} A normalized, three-dimensional host-level hygiene signal covering patch posture, AD exposure state, and telemetry freshness, derived from the HygieneBench synthetic fleet schema and calibrated on held-out seeds.

  \item \textbf{C3 --- Empirical evaluation:} A pre-registered, 25-seed synthetic evaluation demonstrating that HygienePrio achieves approximately 31 pp improvement in Precision@50 over EPSS-only (0.509 vs.\ 0.202; non-overlapping BCa 95\% CIs; synthetic evaluation), with patch posture as the dominant HRS dimension ($\sim$20 pp ablation drop) and AD exposure as a substantial secondary contributor ($\sim$9 pp).

  \item \textbf{C4 --- Failure mode analysis and operationalization boundary:} An explicit identification of conditions under which HygienePrio's advantage is smaller (interaction term equivocal; freshness dimension marginal) and the structural boundary conditions (ground truth circularity, synthetic data, 5-seed calibration) that bound the positive finding. This analysis prevents over-generalization and frames the result as a pre-clinical finding rather than a deployment recommendation.
\end{itemize}

\subsection{Paper Organization}

Section~2 reviews background and related work. Section~3 formalizes the problem. Section~4 describes the synthetic dataset and fleet. Section~5 presents the HygienePrio methodology. Section~6 describes the experimental setup. Section~7 reports results. Section~8 discusses findings, limitations, and connections to prior papers. Section~9 addresses threats to validity. Section~11 concludes.
